% =====================================================================
% Axona — A Programmer's Introduction (≈30-minute talk handout)
% David A. Smith, Axona.net
%
% Compile with:  tectonic Axona-Programmer-Intro.tex
% House style: the Axona explainer's tufte-handout format + listings.
% =====================================================================

\documentclass[nobib]{tufte-handout}

\usepackage[utf8]{inputenc}
\usepackage{xcolor}
\usepackage{booktabs}
\usepackage{microtype}
\usepackage{fancyhdr}
\usepackage{etoolbox}
\usepackage{needspace}
\usepackage{listings}

\pretocmd{\section}{\needspace{6\baselineskip}}{}{}
\pretocmd{\subsection}{\needspace{4\baselineskip}}{}{}

\definecolor{rust}{RGB}{185,78,54}
\definecolor{inkmuted}{RGB}{102,102,102}
\definecolor{codebg}{RGB}{248,247,242}

\lstdefinestyle{axonacode}{
  basicstyle=\ttfamily\fontsize{7.6}{9}\selectfont,
  breaklines=true, breakatwhitespace=false, columns=fullflexible,
  keepspaces=true, showstringspaces=false, frame=single,
  rulecolor=\color{inkmuted!40}, backgroundcolor=\color{codebg},
  framexleftmargin=4pt, framexrightmargin=4pt,
  xleftmargin=2pt, aboveskip=0.6em, belowskip=0.6em,
  commentstyle=\color{inkmuted}\itshape, keywordstyle=\color{rust},
}
\lstset{style=axonacode}

% Roman-numeral sections, to read as a talk agenda.
\setcounter{secnumdepth}{1}
\renewcommand{\thesection}{\Roman{section}}

\begin{document}

\fancypagestyle{axonaplain}{%
  \fancyhf{}\renewcommand{\headrulewidth}{0pt}\fancyfoot[C]{\thepage}%
}
\pagestyle{axonaplain}

\begin{center}
{\fontsize{34}{38}\selectfont \textsc{Axona}\par}
\vspace{0.35em}
{\Large\itshape A Programmer's Introduction\par}
\vspace{0.5em}
{\small David A.~Smith \quad\textperiodcentered\quad \emph{Axona.net} \quad\textperiodcentered\quad v0.1 \quad\textperiodcentered\quad a $\sim$30-minute talk\par}
\end{center}

\vspace{0.3em}
\noindent\begin{quote}\itshape\color{inkmuted}
A 30-minute tour for developers: what Axona is, how it works, what protects it,
the API you actually call, and a live app --- \textsc{Axona Minimal} --- built in
about sixty lines. Follow along at \texttt{demo.axona.net}.
\end{quote}
\vspace{0.4em}

% =====================================================================
\section{What Axona is}

Axona is a peer-to-peer substrate for \textbf{addressing, routing, and
publish/subscribe} that runs with no platform in the middle --- including in an
ordinary web browser.\sidenote{Pure JS/WASM. The reference peer is a web page;
the same kernel runs in Node and in the simulator.} You get a global pub/sub bus
where any peer can publish to a topic and any peer can subscribe to it, and the
messages flow peer-to-peer over an encrypted mesh.

Three things make it different from a hosted message broker:

\begin{itemize}
\item \textbf{Self-authenticating.} A peer's identity \emph{is} its
cryptographic keypair; its node address is derived from its public key and
verified directly by other peers. There is no account, no certificate authority,
no registrar.\sidenote{You are your key. Nothing to sign up for, nothing to
revoke through.}
\item \textbf{Geo-aware, not geo-fenced.} The top of every address is a coarse
geographic cell, so routing favors locality --- an ``area code,'' not a wall.
\item \textbf{No central chokepoint.} A bootstrap broker introduces new peers,
but the mesh carries the traffic itself and converges its own state. The first
version is live; anyone can connect.
\end{itemize}

% =====================================================================
\section{How it works}

\textbf{Addresses.} Every node has a 264-bit id. The \emph{top 8 bits} are an S2
geographic cell (one of 192 covering the planet); the rest is
\textsc{sha-256} of the node's public key.\sidenote{So the address is
\texttt{geocell\,$\|$\,hash(pubkey)} --- location-biased at the top, key-bound
below. Topic ids are built the same way (see \S V).} Closeness is XOR distance,
which the top byte dominates --- so ``nearby in the keyspace'' tracks ``nearby on
Earth,'' \emph{softly}: it biases routing without partitioning anyone out.

\textbf{The mesh.} In the browser, peers form a WebRTC mesh. A WebSocket
\emph{bridge} does the initial introduction (and relays signaling), but it is a
bootstrap rendezvous, not a router --- peers can even open new links through
existing peers without it. Each node keeps a learned routing table (its
\emph{synaptome}) and reaches any id in a few hops.

\textbf{Pub/sub.} A topic maps to an id; the \textbf{R} nodes closest to that id
are its \emph{root set} (replicas). A publish goes to the roots, which fan it out
down a tree of sub-axons to the subscribers.\sidenote{Replication + a fan-out
tree is what makes delivery survive churn: lose a root or a branch and the others
still carry the message.} Three mechanisms keep replicas consistent --- gap-safe
replay on (re)subscribe, root-to-root anti-entropy, and kill convergence --- so a
message is backfilled rather than lost, and a retraction stays retracted.

% =====================================================================
\section{The security model}

Everything below is enforced by cryptography the peers carry themselves --- no
trust server sits behind any of it.

\begin{itemize}
\item \textbf{Identity = keypair.} The node id is derived from the public key, so
a peer cannot claim to be another; the transport handshake is mutually
authenticated and channel-bound (a man-in-the-middle that swaps the media
fingerprint fails the bind).
\item \textbf{Signed publishes, verified at ingress.} A signed message is
checked against the publisher's key at the \emph{root}, before it is cached or
fanned out --- spoofed-signature spam dies at the edge.
\item \textbf{You can only act for yourself.} A subscribe/unsubscribe is honored
only for the authenticated channel peer's own id, so no one can subscribe a
victim (amplification) or unsubscribe them (silencing).
\item \textbf{Bounded hosting.} A node only becomes a root for topics it is
genuinely closest to, and inbound size/count caps apply --- a flood of junk
topics can't make every node allocate unbounded state.
\item \textbf{Content-addressed + fresh.} The message id is the hash of its
content; a per-publisher sequence and TTL drop stale or replayed envelopes.
\item \textbf{Sybil / placement defense.} A memory-hard proof-of-work prices the
\emph{publish identity}, decoupled from the node address so it can't be used to
grind a favorable position in the keyspace.
\end{itemize}

\noindent\textbf{What it does \emph{not} promise.} A kill is best-effort
redaction, not a cryptographic un-send --- a peer that already holds the bytes
keeps them. And a medium no platform can censor is, by the same token, one no
platform can moderate. Design your app with that in mind.

% =====================================================================
\section{The API you'll actually use}

The whole surface is small. You construct a peer, then publish and subscribe.

\smallskip
\noindent\begin{tabular}{@{}ll@{}}
\toprule
\texttt{deriveIdentity(\{lat,lng\})} & make a keypair + 264-bit id \\
\texttt{webTransport(\{bridgeUrl, identity\})} & browser transport (mesh + bridge) \\
\texttt{new AxonaPeer(\{domain,node,identity,transport\})} & the peer \\
\texttt{peer.start()} / \texttt{peer.leave()} & join / drain + leave \\
\midrule
\texttt{peer.pub(topic, msg, \{publisher\})} & publish $\rightarrow$ \texttt{msgId} \\
\texttt{peer.sub(topic, handler, \{publisher,since\})} & subscribe $\rightarrow$ \texttt{Subscription} \\
\texttt{peer.unsub(topic, \{publisher\})} & leave a topic \\
\texttt{peer.pull(msgId\,|\,null, \{topic,publisher\})} & fetch by id, or latest \\
\texttt{peer.kill(topic, msgId, \{publisher\})} & retract (creator-only) \\
\texttt{peer.host(topic?)} & serve a topic without consuming it \\
\texttt{peer.metrics / peer.health()} & counters / introspection \\
\bottomrule
\end{tabular}

\smallskip
\textbf{Where a topic lives} is set by \texttt{opts.publisher}: omit it to
publish under \emph{your own} feed, pass \texttt{null} for a well-known public
topic, or pass a 66-char hex publisher to read/write \emph{someone else's} feed.
A shared ``room'' uses a region-anchored synthetic publisher so every peer in
that region derives the same topic id.\sidenote{The helper
\texttt{resolveAnchor()} returns that publisher from \texttt{?region=} (default
us-east), along with the lat/lng to seed your identity.}

\textbf{The envelope} a subscriber receives is \texttt{\{ msgId, seq, ts, topic,
message, publisher, signerPubkey \}}; a retraction arrives on the same handler as
\texttt{\{ deleted: true, msgId \}}.

% =====================================================================
\section{Build-along: \textsc{Axona Minimal}}

A topic field, a message field, a received-messages pane. Live at
\texttt{demo.axona.net/apps/axona-minimal/} --- open it in two tabs and watch a
message cross between them. Here is the whole of it.

\textbf{1. Connect} --- derive an identity, open the transport, build the peer,
wait briefly for the mesh:

\begin{fullwidth}
\begin{lstlisting}
import { AxonaPeer, AxonaDomain, NeuronNode, deriveIdentity } from '/src/index.js';
import { webTransport } from '/src/transport/web/index.js';
import { resolveAnchor } from '../lib/region.js';

const BRIDGE = 'wss://bridge.axona.net';
const ANCHOR = resolveAnchor();                       // { name, center:{lat,lng}, publisher }

const identity  = await deriveIdentity({ lat: ANCHOR.center.lat, lng: ANCHOR.center.lng });
const transport = webTransport({ bridgeUrl: BRIDGE, identity });
const node      = new NeuronNode({ id: BigInt('0x' + identity.id), lat: ANCHOR.center.lat, lng: ANCHOR.center.lng });
node.transport  = transport;
const peer = new AxonaPeer({ domain: new AxonaDomain({ k: 20 }), node, identity, transport });

await transport.start(identity.id);
await peer.start();                                   // mesh forms over the next moment
\end{lstlisting}
\end{fullwidth}

\textbf{2. Subscribe} to the topic. Re-subscribing to a new topic drops the old
one; the handler gets each envelope:

\begin{fullwidth}
\begin{lstlisting}
currentSub = await peer.sub(topic, (env) => {
  if (!env || env.deleted || seen.has(env.msgId)) return;   // skip our own echo
  seen.add(env.msgId);
  render(env.message, env.signerPubkey, /*self=*/false);
}, { publisher: ANCHOR.publisher, since: 'all' });
\end{lstlisting}
\end{fullwidth}

\textbf{3. Publish.} Own publishes may not echo back, so render optimistically
and let the \texttt{seen} set dedup if they do:

\begin{fullwidth}
\begin{lstlisting}
const msgId = await peer.pub(topic, text, { publisher: ANCHOR.publisher });
seen.add(msgId);
render(text, null, /*self=*/true);
\end{lstlisting}
\end{fullwidth}

That is the entire networking story: \texttt{deriveIdentity} $\rightarrow$
\texttt{webTransport} $\rightarrow$ \texttt{AxonaPeer} $\rightarrow$
\texttt{sub}/\texttt{pub}. Everything else in the file is DOM glue.

% =====================================================================
\section{Try it, and go deeper}

\begin{itemize}
\item \textbf{Run it:} \texttt{demo.axona.net/apps/axona-minimal/} (and the other
demos on that page). Add \texttt{?region=uswest} to land in another keyspace.
\item \textbf{Read it:} the \emph{Explainer} (how the routing works), the
\emph{Architecture} note (kernel $\cdot$ transport $\cdot$ bridge $\cdot$ wire),
and the \emph{API Reference} for every exported symbol.
\item \textbf{Host it:} run a relay (\texttt{axona-relay}) to carry topics in your
region, or stand up your own bridge --- the only piece you need to bootstrap.
\end{itemize}

\vspace{0.3em}
\noindent{\color{rust}\bfseries Sixty lines and a browser tab put you on the
network. There is no step zero.}

\end{document}
